\section{Zadanie 9}
 Niech $\nu : \mathcal{B}(\mathbb{R}) \to \{0,1\}$ będzie zdefiniowana następująco. 
Dla $B \in \mathcal{B}(\mathbb{R})$ mamy $\nu(B) = 1$ wtedy i tylko wtedy, gdy istnieje $\varepsilon > 0$ taki, że $(0,\varepsilon) \subseteq B$. 
Czy funkcja $\nu$ jest skończenie addytywna? 
Jeśli tak, to czy jest przeliczalnie addytywna?
\subsection*{Rozwiązanie}

Niech $E_1 = \mathbb{Q} \cap (0,1) \in \mathcal{B}(\mathbb{R})$ i $E_2 = (0,1) \setminus \mathbb{Q} \in \mathcal{B}(\mathbb{R})$.

$\nu(E_1) = \nu({E_2}) = 0$, a $\nu(E_1 \cup E_2) = \nu( (0,1) ) = 1 \neq \nu(E_1) + \nu(E_2) = 0$.

\subsection*{Odpowiedź}
\[
\boxed{\text{Nie jest skończenie addytywna}}
\]
